\uuid{RBZ5}
\exo7id{1123}
\titre{exo7 1123}
\auteur{liousse}
\organisation{exo7}
\datecreate{2003-10-01}
\isIndication{false}
\isCorrection{false}
\chapitre{Déterminant, système linéaire}
\sousChapitre{Calcul de déterminants}
\cpge{mpsi-mp2i.determinants.b, pcsi.determinants.a, ptsi.determinants.c, mpsi-mp2i.determinants.d, pcsi.calcul-matriciel-systemes-lineaires.d, ptsi.calcul-matriciel-systemes-lineaires.d, mpsi-mp2i.determinants.e, pcsi.determinants.d, bcpst1.matrices, pt.determinants.a, bcpst2.applications-lineaires-matrices.c}
\module{Algèbre}
\niveau{L2}
\difficulte{}

\contenu{
\question{
Soit $n$ un entier sup\'erieur ou 
\'egal \`a 3. On se place dans $\Rr^{n}.$ On note 
$e_i$ le vecteur de $\Rr^{n}$ dont la i-i\`eme  composante est 
\'egale \`a 1 et toutes les autres sont nulles. \'Ecrire la matrice $n\times n$ 
dont les vecteurs colonnes $C_i$ sont donn\'es par 
$C_i= e_i+e_n$ pour $1\leq i\leq n-1$ et $C_n=e_1+e_2+e_n.$ 
Calculer alors son d\'eterminant.
}
}
